\documentclass[10pt,a4paper]{amsart}
\usepackage[margin=19mm]{geometry}
\usepackage{amsmath,amssymb,mathtools}
\usepackage[scr=rsfs,cal=euler]{mathalfa}
\usepackage{xfrac}
\usepackage[unicode,hidelinks,bookmarks=false]{hyperref}
\newcommand{\ie}{i.e.\ }
\newcommand{\cf}{cf.\ }
\newcommand{\resp}{respectively\ }
\newcommand{\Subsection}{\subsection}
\providecommand{\nobreakdash}{\nobreak\hbox{-}}
\setlength{\emergencystretch}{2em}
\multlinegap0pt
\usepackage{tikz}
\usetikzlibrary{cd}
\usepackage{amssymb}
\usepackage{dsfont}
\usepackage{xfrac}
\usepackage{mathtools}
\usepackage{tensor}
\usepackage{tikz}
\usepackage{xcolor}
\newtheorem{thm}{Theorem}
\newcommand{\m}[1]{{\mathrm{#1}}}
\title{Spin(7)-Orbifold Resolutions}
\author{Viktor F. Majewski}
\date{Source: arXiv:2509.16057v3 (CC BY 4.0)}
\begin{document}
\maketitle

\noindent\emph{Source and licence.} Spin(7)-Orbifold Resolutions. Version arXiv:2509.16057v3 (CC BY 4.0), distributed by arXiv under the Creative Commons Attribution 4.0 International licence, http://creativecommons.org/licenses/by/4.0/, as recorded in the arXiv licence metadata for this exact version and shown on the abstract page https://arxiv.org/abs/2509.16057v3. Full licence notice: \texttt{licenses/arxiv-2509.16057-CC-BY-4.0.txt}.\par\smallskip

\noindent\textit{Editorial scope.} Counted unit: the article's thm codimensionfourMcKayDuality (source0.tex lines 3845--3852). The article's complete original proof of it is delivered with it (source0.tex lines 3854--3985, one \texttt{proof} environment of 132 lines). No mathematics is written, repaired or re-derived: the statement and the proof are the article's own lines, in the article's own notation, and no retained line is substituted or re-flowed.  No further source-local context is needed: every cross-reference of every retained line resolves inside the two delivered spans.  Nothing was omitted from any retained span, so the package carries no omission record.  No retained line contains a citation, so the wrapper carries no bibliography and no bibliography entry.  The retained lines contain the article's own diagram code, which the wrapper typesets with the standard tikz packages; no image file is delivered and the package declares no asset dependency.  Editorial formatting only, in addition: an AMS \texttt{amsart} preamble that restates the article's own declarations and macros used by the retained lines, namely the environments \texttt{thm} on separate counters, together with the standard packages those lines need.  The retained environments print in this document as Theorem 1 in order of appearance, whereas the article prints the same items with the numbering of its own full document.  The complete native source of the article as distributed in the arXiv e-print for this exact version is bundled unchanged as \texttt{sources/185346/source0.tex}.
\begin{thm}
\label{codimensionfourMcKayDuality}
    Let $[\zeta]\in \m{H}^\bullet(S,\mathfrak{H}^\bullet_\Gamma)$ corresponds to a $g_S$-harmonic $\mathfrak{H}_\Gamma$-twisted differential form $\zeta\in \Gamma(S,\mathfrak{P}_\Gamma)\subset \Omega^\bullet(S,\mathfrak{H}^\bullet_\Gamma)$. Then the ACF-space 
    \begin{align*}
        \rho_\zeta\colon (N_\zeta,\Phi_\zeta)\dashrightarrow (N_0,\Phi_0)
    \end{align*}
    is adiabatic torsion-free. Conversely, $\rho_\zeta\colon (N_\zeta,\Phi_\zeta)\dashrightarrow (N_0,\Phi_0)$ is adiabatic torsion-free if $\zeta\in \Omega^\bullet(S,\mathfrak{H}^\bullet_\Gamma)$ is harmonic and hence determines a unique class $[\Phi]\oplus[\zeta]\in \m{H}^4_{\m{CR}}(X,\mathbb{R})\cong \m{H}^4_{\m{dR}}(X,\mathbb{R})\oplus \m{H}^2(S,\mathfrak{H}^2_\Gamma)$.
\end{thm}

\begin{proof}
We define $\overline{\zeta}\in\Omega^0(\m{Fr}(X/S,\Phi),\Theta_{\m{Im}(\mathbb{H})}(\Gamma))^{\m{N}_{\m{Spin}(7)}(\Gamma)}$ via 
\begin{align*}
    \zeta=\left<\overline{\zeta},\left[\theta_{S}\wedge\theta_S\right]_+\right>
\end{align*}
and the identification of $\Theta_{\m{Im}(\mathbb{H})}(\Gamma)=\m{Im}(\mathbb{H})\otimes\Theta_{\m{Im}(\mathbb{C})}(\Gamma)\cong \m{Im}(\mathbb{H})\otimes \mathfrak{z}(\mathfrak{pu}(R_\Gamma)^{\Gamma})$.\\

We will lift the forms from $N_\zeta$ to the bundle $\overline{\zeta}^*\mathcal{M}_\Gamma$, where $\overline{\zeta}\colon \m{Fr}(X/S,\Phi)\rightarrow \Theta_{\m{Im}(\mathbb{H})}$ denote the $\m{N}_{\m{Spin}(7)}(\Gamma)$-equivariant map lifting $\zeta$.


Now, consider the pullback of $\kappa \colon\mathcal{M}_\Gamma\rightarrow \Theta_{\m{Im}(\mathbb{H})}$ along $\overline{\zeta}$.
    \begin{equation*}
         \begin{tikzcd}
            \overline{\zeta}^*\mathcal{M}_\Gamma\arrow[r,"\kappa^*\overline{\zeta}"]\arrow[d,swap,"\overline{\zeta}^*\kappa"]&\mathcal{M}_\Gamma\arrow[d,"\kappa"]\\
            \m{Fr}(X/S,\Phi)\arrow[r,"\overline{\zeta}"]&\Theta_{\m{Im}(\mathbb{H})}(\Gamma)
        \end{tikzcd}
    \end{equation*}

As the group $\m{N}_{\m{Spin}(7)}(\Gamma)$ acts freely from the right on $\m{Fr}(X/S,\Phi)\times \mathcal{M}_\Gamma$, it acts freely on $\overline{\zeta}^*\mathcal{M}_\Gamma$ and the projection $\psi_\zeta \colon\overline{\zeta}^*\mathcal{M}_\Gamma\rightarrow N_\zeta$ defines a principal $\m{N}_{\m{Spin}(7)}(\Gamma)$-bundle.

\begin{equation*}
    \begin{tikzcd}
        &\overline{\zeta}^*\mathcal{M}_\Gamma\arrow[dl,swap,"\psi_\zeta"]\arrow[r,"\kappa^*\overline{\zeta}"]\arrow[d,"\overline{\zeta}^*\kappa"]&\mathcal{M}_\Gamma\arrow[d,"\kappa"]\\
        N_\zeta\arrow[dr,"\nu_\zeta"]&\m{Fr}(X/S,\Phi)\arrow[r,"\overline{\zeta}"]\arrow[d,"\phi"]&\Theta_{\m{Im}(\mathbb{H})}(\Gamma)\\
        &S&
    \end{tikzcd}
\end{equation*}
where 
\begin{align*}
    \nu_\zeta\colon N_\zeta=\overline{\zeta}^*\mathcal{M}_\Gamma/\m{N}_{\m{Spin}(7)}(\Gamma)\rightarrow S
\end{align*}
denotes the quotient mapping onto $S$.\\

The Levi-Civita connection on $S$ pulls back to a connection $(\overline{\zeta}^*\kappa)^*\varphi=\varphi_\zeta$ on $\psi_\zeta \colon\overline{\zeta}^*\mathcal{M}_\Gamma\rightarrow N_\zeta$. Furthermore, 
the $\m{N}_{\m{Spin}(7)}(\Gamma)$-invariant connection ${\check{H}}$ on $\kappa \colon\mathcal{M}_\Gamma\rightarrow \Theta_{\m{Im}(\mathbb{H})}(\Gamma)$ gives rise to a splitting of the horizontal subbundle of $(\overline{\zeta}^*\kappa)^*\varphi$. Consequently, the tangent bundle of $\overline{\zeta}^*\mathcal{M}_\Gamma$ decomposes into 
\begin{align*}
    T\overline{\zeta}^*\mathcal{M}_\Gamma\cong&(\overline{\zeta}^*\kappa)^*H_S\oplus\overline{\zeta}^*K\oplus\underline{\mathfrak{n}_{\m{Spin}(7)}(\Gamma)}_{\overline{\zeta}^*\mathcal{M}_\Gamma}\\
    \eqqcolon&T^{1,0,0}\overline{\zeta}^*\mathcal{M}_\Gamma\oplus T^{0,1,0}\overline{\zeta}^*\mathcal{M}_\Gamma\oplus T^{0,0,1}\overline{\zeta}^*\mathcal{M}_\Gamma.
\end{align*}
In particular, $(\overline{\zeta}^*\kappa)^*H_S$ defines an $\m{N}_{\m{Spin}(7)}(\Gamma)$ equivariant connection on $\phi\circ\overline{\zeta}^*\kappa \colon\overline{\zeta}^*\mathcal{M}_\Gamma\rightarrow S$, and thus descends to an Ehresmann connection $H_\zeta$ on $\nu_\zeta\colon N_\zeta\rightarrow S$.\\

Differential forms on $\overline{\zeta}^*\mathcal{M}_\Gamma$ decompose into
\begin{align*}
    \Omega^{p,q,r}(\overline{\zeta}^*\mathcal{M}_\Gamma)=\Gamma(\overline{\zeta}^*\mathcal{M}_\Gamma,\wedge^p(\overline{\zeta}^*\kappa)^* H_S^* \otimes \wedge^p(\kappa^*\overline{\zeta})^*K^* \otimes\wedge^r\underline{\mathfrak{n}_{\m{Spin}(7)}(\Gamma)}_{\overline{\zeta}^*\mathcal{M}_\Gamma}^*)
\end{align*}
and the de Rham differential can be identified with 
\begin{equation*}
    \m{d}
    =\begin{cases}
        \m{d}^{1,0}&=\m{d}^{1,0,0}+\m{d}^{0,1,0}+\m{d}^{2,-1,0}\\
        \m{d}^{0,1}&=\m{d}^{0,0,1}\\
        \m{d}^{2,-1}&=\m{d}^{2,0,-1}.
    \end{cases}
\end{equation*}


We identify $\Phi_\zeta$ with a four form on $\overline{\zeta}^*\mathcal{M}_\Gamma$ given by 
\begin{align*}
    \hat{\Phi}_\zeta^{4,0}&=(\phi\circ(\overline{\zeta}^*\kappa))^*\m{vol}_S\\
    \hat{\Phi}_\zeta^{2,2}&=-\left<(\overline{\zeta}^*\kappa)^*\left[\theta_{S}\wedge\theta_S\right]_+\wedge((\kappa^*\overline{\zeta})^*\underline{\omega}_{\mathcal{M}_\Gamma})^{0,2,0}\right>\\
    \hat{\Phi}_\zeta^{0,4}&=\frac{1}{6}\left<((\kappa^*\overline{\zeta})^*\underline{\omega}_{\mathcal{M}_\Gamma})^{0,2,0}\wedge((\kappa^*\overline{\zeta})^*\underline{\omega}_{\mathcal{M}_\Gamma})^{0,2,0}\right>.
\end{align*}
We will now proceed by computing the torsion $\m{d}\Phi_\zeta$ and prove that if $\zeta$ is harmonic, $\m{d}\Phi_\zeta=\m{d}^{2,-1}\Phi^{2,2}_\zeta+\m{d}^{2,-1}\Phi^{0,4}_\zeta$.

First we compute
\begin{align*}
    \psi_\zeta^*\m{d}\Phi_\zeta^{4,0}=&\m{d}(\overline{\zeta}^*\kappa\circ\phi)^*\m{vol}_S\\
    =&\m{d}(\overline{\zeta}^*\kappa\circ\phi)^*\m{vol}_S\\
    =&(\overline{\zeta}^*\kappa\circ\phi)^*\m{d}\m{vol}_S\\
    =&0
\end{align*}
To see that the harmonicity of $\zeta$ implies that $\Phi_\zeta$ defines an adiabatic torsion-free $\m{Spin}(7)$-structure, we compute

\begin{align*}
    \m{d}_{\varphi}\zeta=&\m{d}^{1,0}\left<\left[\theta_{S}\wedge\theta_S\right]_+,\overline{\zeta}\right>=\left<\left[\theta_{S}\wedge\theta_S\right]_+\wedge(\overline{\zeta}^*\theta_{\Theta_{\m{Im}(\mathbb{H})}(\Gamma)})^{1,0}\right>.
\end{align*}
where $\theta_{\Theta_{\m{Im}(\mathbb{H})}(\Gamma)}\in \Omega^1(\Theta_{\m{Im}(\mathbb{H})}(\Gamma),\m{Im}(\mathbb{H})\otimes \mathfrak{z}(\mathfrak{pu}(R_\Gamma)^{\Gamma}))$ denotes the Soldering form on $\Theta_{\m{Im}(\mathbb{H})}(\Gamma)$. We use that $\varphi$ is torsion-free, as well as
\begin{align*}
    \m{d}^{1,0,0}((\kappa^*\overline{\zeta})^*\underline{\omega}_{\mathcal{M}_\Gamma})^{0,2,0}=&(\m{d}(\kappa^*\overline{\zeta})^*\underline{\omega}_{\mathcal{M}_\Gamma})^{1,2,0}-\m{d}^{0,1,0}((\kappa^*\overline{\zeta})^*\underline{\omega}_{\mathcal{M}_\Gamma})^{1,1,0}\\
    =&((\kappa^*\overline{\zeta})^*\m{d}\underline{\omega}_{\mathcal{M}_\Gamma})^{1,2,0}\\
    &-\m{d}^{0,1,0}((\kappa^*\overline{\zeta})^*\underline{\omega}_{\mathcal{M}_\Gamma})^{1,1,0}\\
    =&-\left((\kappa^*\overline{\zeta})^*\left<\kappa^*\theta_{\Theta_{\m{Im}(\mathbb{H})}(\Gamma)}\wedge\pi_{\mathfrak{z}(\mathfrak{pu}(R_\Gamma)^{\Gamma})}F_{\mathfrak{a}_\Gamma}\right>\right)^{1,2,0}\\
    &-\m{d}^{0,1,0}((\kappa^*\overline{\zeta})^*\underline{\omega}_{\mathcal{M}_\Gamma})^{1,1,0}\\
    =&-\left<\overline{\zeta}^*(\kappa^*\overline{\zeta})^*\theta_{\Theta_{\m{Im}(\mathbb{H})}(\Gamma)}\wedge(\kappa^*\overline{\zeta})^*\pi_{\mathfrak{z}(\mathfrak{pu}(R_\Gamma)^{\Gamma})}F_{\mathfrak{a}_\Gamma}\right>^{1,2,0}\\
    &-\m{d}^{0,1,0}((\kappa^*\overline{\zeta})^*\underline{\omega}_{\mathcal{M}_\Gamma})^{1,1,0}\\
    =&-\left<(\overline{\zeta}^*(\kappa^*\overline{\zeta})^*\theta_{\Theta_{\m{Im}(\mathbb{H})}(\Gamma)})^{1,0,0}\wedge((\kappa^*\overline{\zeta})^*\pi_{\mathfrak{z}(\mathfrak{pu}(R_\Gamma)^{\Gamma})}F_{\mathfrak{a}_\Gamma})^{0,2,0}\right>\\
    &-\left<(\overline{\zeta}^*(\kappa^*\overline{\zeta})^*\theta_{\Theta_{\m{Im}(\mathbb{H})}(\Gamma)})^{0,1,0}\wedge((\kappa^*\overline{\zeta})^*\pi_{\mathfrak{z}(\mathfrak{pu}(R_\Gamma)^{\Gamma})}F_{\mathfrak{a}_\Gamma})^{1,1,0}\right>\\
    &-\m{d}^{0,1,0}((\kappa^*\overline{\zeta})^*\underline{\omega}_{\mathcal{M}_\Gamma})^{1,1,0}
\end{align*}
whereby 
\begin{align*}
    \m{d}^{0,1,0}((\kappa^*\overline{\zeta})^*\underline{\omega}_{\mathcal{M}_\Gamma})^{1,1,0}=-\left<(\overline{\zeta}^*(\kappa^*\overline{\zeta})^*\theta_{\Theta_{\m{Im}(\mathbb{H})}(\Gamma)})^{0,1,0}\wedge((\kappa^*\overline{\zeta})^*\pi_{\mathfrak{z}(\mathfrak{pu}(R_\Gamma)^{\Gamma})}F_{\mathfrak{a}_\Gamma})^{1,1,0}\right>
\end{align*}

and 
\begin{align*}
    \m{d}^{0,1,0}((\kappa^*\overline{\zeta})^*\underline{\omega}_{\mathcal{M}_\Gamma})^{0,2,0}=&(\m{d}(\kappa^*\overline{\zeta})^*\underline{\omega}_{\mathcal{M}_\Gamma})^{0,3,0}\\
    =&-\left<(\kappa^*\overline{\zeta})^*\left(\kappa^*\theta_{\Theta_{\m{Im}(\mathbb{H})}(\Gamma)}\wedge\pi_{\mathfrak{z}(\mathfrak{pu}(R_\Gamma)^{\Gamma})}F_{\mathfrak{a}_\Gamma}\right)\right>^{0,3,0}\\
    =&-\left<(\overline{\zeta}^*(\kappa^*\overline{\zeta})^*\theta_{\Theta_{\m{Im}(\mathbb{H})}(\Gamma)})^{0,1,0}\wedge((\kappa^*\overline{\zeta})^*\pi_{\mathfrak{z}(\mathfrak{pu}(R_\Gamma)^{\Gamma})}F_{\mathfrak{a}_\Gamma})^{0,2,0}\right>.
\end{align*}
Hence, we deduce that 

\begin{align*}
    (\m{d}^{1,0,0}+\m{d}^{0,1,0})((\kappa^*\overline{\zeta})^*\underline{\omega}_{\mathcal{M}_\Gamma})^{0,2,0}=&\left<(\overline{\zeta}^*\kappa)^*(\overline{\zeta}^*\theta_{\Theta_{\m{Im}(\mathbb{H})}(\Gamma)})^{1,0}\wedge((\kappa^*\overline{\zeta})^*\pi_{\mathfrak{z}(\mathfrak{pu}(R_\Gamma)^{\Gamma})}F_{\mathfrak{a}_\Gamma})^{0,2,0}\right>.
\end{align*}

Further, as the pull back of the Soldering form satisfies
\begin{align*}
    \m{d}(\overline{\zeta}^*\kappa)^*\left[\theta_{S}\wedge\theta_S\right]_+=&(\overline{\zeta}^*\kappa)^*\m{d}\left[\theta_{S}\wedge\theta_S\right]_+\\
    =&(\overline{\zeta}^*\kappa)^*[\varphi\overset{\circ}{\wedge}\left[\theta_{S}\wedge\theta_S\right]_+]\in\Omega^{2,1}(\overline{\zeta}^*\mathcal{M}_\Gamma,\m{Im}(\mathbb{H}))^{\m{N}_{\m{Spin}(7)}(\Gamma)}
\end{align*}
we deduce, that if $\zeta$ is harmonic
\begin{align*}
    \psi_\zeta^*\m{d}\Phi^{2,2}_\zeta=&-\left<(\overline{\zeta}^*\kappa)^*\left[\theta_{S}\wedge\theta_S\right]_+\wedge\left<(\overline{\zeta}^*\kappa)^*(\overline{\zeta}^*\theta_{\Theta_{\m{Im}(\mathbb{H})}(\Gamma)})^{1,0}\wedge((\kappa^*\overline{\zeta})^*\pi_{\mathfrak{z}(\mathfrak{pu}(R_\Gamma)^{\Gamma})}F_{\mathfrak{a}_\Gamma})^{0,2,0}\right>\right>\\
    &+\psi^*_\zeta\m{d}^{2,-1}\Phi^{2,2}_\zeta\\
    =&-\left<(\overline{\zeta}^*\kappa)^*\underbrace{\left<\left[\theta_{S}\wedge\theta_S\right]_+\wedge(\overline{\zeta}^*\theta_{\Theta_{\m{Im}(\mathbb{H})}(\Gamma)})^{1,0}\right>}_{\m{d}_{\varphi}\zeta=0}\wedge((\kappa^*\overline{\zeta})^*\pi_{\mathfrak{z}(\mathfrak{pu}(R_\Gamma)^{\Gamma})}F_{\mathfrak{a}_\Gamma})^{0,2,0}\right>\\
    &+\psi^*_\zeta\m{d}^{2,-1}\Phi^{2,2}_\zeta\\
    =&\psi^*_\zeta\m{d}^{2,-1}\Phi^{2,2}_\zeta
\end{align*}
Further,
\begin{align*}
    \psi_\zeta^*\m{d}^{1,0}\Phi_\zeta^{0,4}=&\m{d}^{1,0,0}\hat{\Phi}_\zeta^{0,4}\\
    =&\frac{1}{6}\m{d}^{1,0,0}\left<((\kappa^*\overline{\zeta})^*\underline{\omega}_{\mathcal{M}_\Gamma})^{0,2,0}\wedge((\kappa^*\overline{\zeta})^*\underline{\omega}_{\mathcal{M}_\Gamma})^{0,2,0}\right>\\
    =&\frac{1}{3}\left<\m{d}^{1,0,0}((\kappa^*\overline{\zeta})^*\underline{\omega}_{\mathcal{M}_\Gamma})^{0,2,0}\wedge((\kappa^*\overline{\zeta})^*\underline{\omega}_{\mathcal{M}_\Gamma})^{0,2,0}\right>\\
    =&-\frac{1}{3}\left<\left<(\overline{\zeta}^*\kappa)^*(\overline{\zeta}^*\theta_{\Theta_{\m{Im}(\mathbb{H})}(\Gamma)})^{1,0}\wedge((\kappa^*\overline{\zeta})^*\pi_{\mathfrak{z}(\mathfrak{pu}(R_\Gamma)^{\Gamma})}F_{\mathfrak{a}_\Gamma})^{0,2,0}\right>\wedge((\kappa^*\overline{\zeta})^*\underline{\omega}_{\mathcal{M}_\Gamma})^{0,2,0}\right>\\
    =&-\frac{1}{3}\left<(\overline{\zeta}^*\kappa)^*(\overline{\zeta}^*\theta_{\Theta_{\m{Im}(\mathbb{H})}(\Gamma)})^{1,0}\wedge\left<((\kappa^*\overline{\zeta})^*\pi_{\mathfrak{z}(\mathfrak{pu}(R_\Gamma)^{\Gamma})}F_{\mathfrak{a}_\Gamma})^{0,2,0}\wedge((\kappa^*\overline{\zeta})^*\underline{\omega}_{\mathcal{M}_\Gamma})^{0,2,0}\right>\right>\\
    =&-\frac{1}{3}\left<(\overline{\zeta}^*\kappa)^*(\overline{\zeta}^*\theta_{\Theta_{\m{Im}(\mathbb{H})}(\Gamma)})^{1,0}\wedge\left<((\kappa^*\overline{\zeta})^*(\underbrace{\pi_{\mathfrak{z}(\mathfrak{pu}(R_\Gamma)^{\Gamma})}F_{\mathfrak{a}_\Gamma}\wedge\underline{\omega}_{\mathcal{M}_\Gamma}}_{=0}))^{0,4,0}\right>\right>\\
    =&0
\end{align*}

vanishes as a consequence of $F_{\mathfrak{a}_\Gamma}$ being fibrewise ASD.
\end{proof}

\end{document}
